\section{Fundamentals}

\subsection{MOS Transistors}

\begin{definition}
	\textbf{nMOS}: ON when Gate=1, OFF when Gate=0. Passes strong 0, weak 1.
	\textbf{pMOS}: ON when Gate=0, OFF when Gate=1. Passes strong 1, weak 0.
\end{definition}

\begin{definition}
	\textbf{CMOS Logic}: Two complementary networks:
	\begin{itemize}
		\item \textbf{Pull-Up (PUN)}: pMOS connects output to $V_{DD}$ (1).
		\item \textbf{Pull-Down (PDN)}: nMOS connects output to GND (0).
	\end{itemize}
	Output is always strongly driven; no significant static power draw.
\end{definition}

\subsection{Logic Gates}

\begin{definition}
	\textbf{NOT}: $Y = \overline{A}$ (2 transistors).
	\textbf{NAND}: $Y = \overline{A \cdot B}$ (4T: 2 pMOS parallel, 2 nMOS series).
	\textbf{NOR}: $Y = \overline{A + B}$ (4T: 2 pMOS series, 2 nMOS parallel).
	AND/OR are NAND/NOR + inverter (inverting gates are natural in CMOS).
\end{definition}

\begin{remark}
	\textbf{XOR}: $Y = A \oplus B$ (1 if different).
	\textbf{XNOR}: $Y = \overline{A \oplus B}$ (1 if same).
	\textbf{Buffer}: $Y = A$ (two NOTs in series, for amplification/delay).
\end{remark}

\subsection{Boolean Algebra}

\begin{definition}
	Key rules: Identity ($A \cdot 1 = A$, $A + 0 = A$), Complement ($A \cdot \overline{A} = 0$, $A + \overline{A} = 1$), Distributive ($A(B+C) = AB + AC$).
\end{definition}

\begin{theorem}
	\textbf{DeMorgan's}: $\overline{A \cdot B} = \overline{A} + \overline{B}$ and $\overline{A + B} = \overline{A} \cdot \overline{B}$.
\end{theorem}

\begin{important}
	A gate set is \textbf{logically complete} if any Boolean function can be built from it. Complete sets: $\{$AND, OR, NOT$\}$, $\{$NAND$\}$, $\{$NOR$\}$.
\end{important}

\begin{definition}
	\textbf{Minterm}: Product term with all variables. \textbf{SOP}: OR of minterms where output=1. \textbf{K-Map}: Group $2^n$ adjacent 1s to eliminate changing variables.
\end{definition}

\subsection{Combinational Logic}

\begin{definition}
	Output is a pure function of current inputs (no memory/state).
\end{definition}

\begin{definition}
	\textbf{Decoder} ($n$-to-$2^n$): Sets exactly one of $2^n$ outputs to 1 based on $n$-bit input.
	\textbf{Multiplexer}: $\log_2 N$ select lines choose which of $N$ data inputs reaches output. Can implement any logic function as a LUT.
\end{definition}

\begin{definition}
	\textbf{Full Adder}: $S = A \oplus B \oplus C_{in}$, $C_{out} = AB + C_{in}(A \oplus B)$.
	\textbf{Ripple Carry Adder}: Chain $n$ full adders; critical path is carry ripple from LSB to MSB.
\end{definition}

\begin{lemma}
	\textbf{Tri-state buffer}: Output can be 0, 1, or High-Z (disconnected). Enables shared buses.
	\textbf{PLA}: Programmable AND-plane + OR-plane = physical SOP implementation.
\end{lemma}

\subsection{Sequential Logic}

\begin{definition}
	Output depends on current inputs \textbf{and past state}. Uses storage elements (feedback loops).
\end{definition}

\begin{definition}
	\textbf{R-S Latch}: Cross-coupled NANDs. Set ($S$=0) $\to Q$=1, Reset ($R$=0) $\to Q$=0. $S$=$R$=0 is \textbf{forbidden} (metastability).
	\textbf{D Latch}: Adds Write Enable; when WE=1, $Q$ follows $D$. \textbf{Transparent} while enabled.
\end{definition}

\begin{important}
	\textbf{D Flip-Flop}: Master-slave pair of D latches. \textbf{Edge-triggered}: captures $D$ only at rising clock edge. Solves transparency problem for synchronous design.
\end{important}

\subsubsection{Timing and Clock}

\begin{definition}
	\textbf{Asynchronous}: Transitions on input change (no coordination).
	\textbf{Synchronous}: Transitions on clock edges (predictable, global).
	\textbf{Clock}: Periodic 0/1 signal; period must cover worst-case combinational delay.
\end{definition}

\subsubsection{Finite State Machines}

\begin{definition}
	\textbf{FSM}: States + transitions + inputs + outputs.
	\textbf{Moore}: Output depends only on state.
	\textbf{Mealy}: Output depends on state \textbf{and} inputs.
\end{definition}

\begin{remark}
	FSM implementation: (1) State register (sequential), (2) Next-state logic (combinational), (3) Output logic (combinational). State encoding: Binary (min bits), One-Hot (one bit/state), Output (use outputs as state bits).
\end{remark}

\subsubsection{Memory and Registers}

\begin{definition}
	\textbf{Register}: Group of D flip-flops with common clock/WE.
	\textbf{Memory}: Array indexed by address. Address space = $2^n$ locations. Uses decoder + MUX.
\end{definition}

\subsection{Timing}

\subsubsection{Combinational Timing}

\begin{definition}
	\textbf{Propagation delay ($t_{pd}$)}: Max time from input change to output stable (longest path).
	\textbf{Contamination delay ($t_{cd}$)}: Min time from input change to output starts changing (shortest path).
\end{definition}

\begin{remark}
	\textbf{Glitch}: Multiple output transitions from a single input change due to different path delays. Fixed by adding consensus terms in K-maps.
\end{remark}

\subsubsection{Sequential Timing}

\begin{definition}
	\textbf{Aperture time} $t_a = t_{setup} + t_{hold}$:
	$t_{setup}$: Data must be stable \textbf{before} clock edge.
	$t_{hold}$: Data must stay stable \textbf{after} clock edge.
	Violation $\to$ \textbf{metastability} (output stuck between 0/1).
\end{definition}

\begin{definition}
	$t_{pcq}$: Max time after clock edge until $Q$ stable.
	$t_{ccq}$: Min time after clock edge until $Q$ starts changing.
\end{definition}

\begin{theorem}
	\textbf{Setup constraint}: $T_c \geq t_{pcq} + t_{pd} + t_{setup}$ (determines $f_{max}$).
	\textbf{Hold constraint}: $t_{ccq} + t_{cd} > t_{hold}$ (cannot fix by slowing clock).
\end{theorem}

\begin{important}
	\textbf{Clock skew} ($t_{skew}$) worsens both:
	Setup: $T_c \geq t_{pcq} + t_{pd} + t_{setup} + t_{skew}$.
	Hold: $t_{ccq} + t_{cd} > t_{hold} + t_{skew}$.
\end{important}
